%%═══════════════════════════════════════════════════════════════════
%% Lecture 02 — The Prismatic Bar: A Leading Example in Elasticity
%% Course  : Mechanics of Continua (77606)
%% Lecturer: Prof. Barak Kol
%% Date    : 25 March 2026
%%═══════════════════════════════════════════════════════════════════

\section{Lecture 02 — The Prismatic Bar in Uniaxial Stress}
\label{sec:lec02}

%%────────────────────────────────────────────────────────────────────
\subsection*{Overview and Motivation}
%%────────────────────────────────────────────────────────────────────

This lecture revisits stress from the ground up by studying a single,
concrete object: a \emph{prismatic bar} (a rod of constant cross-section)
under \emph{uniaxial loading}.  The bar is our \textbf{leading example}
for the entire theory of elasticity — simple enough to solve exactly, yet
rich enough to introduce every key concept: stress, strain, Young's
modulus, Poisson's ratio, the stress tensor, and the static equilibrium
equations.

We consider a bar of:
\begin{itemize}
    \item length $L$,
    \item uniform cross-sectional area $A$ (square, circular, or any
          shape — the analysis is identical),
    \item isotropic, linearly elastic material (e.g.\ steel, concrete,
          bone).
\end{itemize}
An \textbf{axial load} $F$ is applied inward at both ends (compression)
or outward at both ends (tension).  Static equilibrium means the two
forces are equal and opposite; the bar as a whole does not accelerate.

%%────────────────────────────────────────────────────────────────────
\subsection{Stress and Strain in 1D}
\label{subsec:lec02:1d}
%%────────────────────────────────────────────────────────────────────

\subsubsection*{Internal force and the concept of stress}

Imagine cutting the bar at an arbitrary cross-section and isolating the
right half.  That half is in equilibrium: the external pull $F$ to the
right must be balanced by an equal and opposite \emph{internal contact
force} $F$ to the left, exerted by the left half through the cut surface.
Because the same argument applies at \emph{any} cross-section, the
internal force is equal to $F$ everywhere along the bar.

The key insight is to normalise by area.  We define the \textbf{stress}
\begin{align}
    \sigma \;=\; \frac{F}{A},
    \label{eq:lec02:stress_def}
\end{align}
which is the force per unit area acting on the cross-section.  Stress
has the same SI units as pressure — \emph{pascals} (Pa, named after
Blaise Pascal, 1647):
\[
    [\sigma] = \frac{\text{N}}{\text{m}^2} = \text{Pa}.
\]
Typical engineering values are in the range of megapascals (MPa) to
gigapascals (GPa).

\textbf{Physical intuition:} Two bars made of the same material but with
different cross-sections, loaded by the same force, deform by the same
\emph{strain} only if acted on by the same \emph{stress}.  It is stress,
not force, that drives deformation.

\subsubsection*{Deformation of a volume element and strain}

Consider an infinitesimal volume element of length $dz$ and area $dA$.
It is pulled by forces $dF = \sigma\,dA$ on both faces.  By the same
proportionality argument for any element,

\begin{itemize}
    \item Each element elongates by some amount $d(\Delta z)$ proportional
          to its length $dz$.
    \item Summing over all elements, the total elongation $\Delta L$ is
          proportional to the original length $L$.
\end{itemize}

We define the \textbf{longitudinal strain} (also called \emph{elongation}) as the dimensionless relative deformation,
\begin{align}
    \varepsilon_{zz} \;\equiv\; \frac{\Delta L}{L}.
    \label{eq:lec02:strain_def}
\end{align}

\textbf{Note on notation:} Landau \& Lifshitz use a different
(older) convention; this course adopts the modern standard in which
strain is denoted $\varepsilon_{ij}$.

\subsubsection*{Hooke's Law and Young's modulus}

Linear elastic response means stress is proportional to strain:
\begin{align}
    \boxed{\sigma \;=\; E\,\varepsilon_{zz}},
    \label{eq:lec02:hooke1d}
\end{align}
where $E$ is \textbf{Young's modulus} (or the \emph{elastic modulus}),
first proposed by Euler (1727) and systematised by Thomas Young (1807).
Re-substituting the definitions of $\sigma$ and $\varepsilon_{zz}$ gives
the fundamental displacement formula:
\begin{align}
    \boxed{\Delta L \;=\; \frac{F\,L}{E\,A}}.
    \label{eq:lec02:elongation}
\end{align}

\textbf{The analogy with resistors:}  
The electrical resistance of a wire is $R = \rho_e\,L/A$, where $\rho_e$
is the resistivity.  Under a voltage $V$, the current is $I = V/R$.
Equation~\eqref{eq:lec02:elongation} has the identical structure with
$F \leftrightarrow V$, $E \leftrightarrow 1/\rho_e$, and
$\Delta L \leftrightarrow I$.  This is not a coincidence — both
phenomena are governed by linear constitutive relations in geometrically
similar media.

\textbf{Dimensions of $E$:}  Since $\varepsilon_{zz}$ is dimensionless,
$E$ has the same dimensions as stress:
\[
    [E] = \text{Pa} = \frac{\text{N}}{\text{m}^2}
        = \frac{\text{J}}{\text{m}^3}.
\]
The last equality shows that $E$ can also be read as energy stored per
unit volume — a useful perspective in thermodynamics of solids.

%%────────────────────────────────────────────────────────────────────
\subsection{Lateral Contraction and Poisson's Ratio}
\label{subsec:lec02:poisson}
%%────────────────────────────────────────────────────────────────────

When a bar is stretched longitudinally, it contracts in the lateral
(transverse) directions — and vice versa.  By symmetry, both transverse
directions $x$ and $y$ contract equally.  The \textbf{Poisson's ratio}
(named after Siméon Denis Poisson, 1827) quantifies this coupling:
\begin{align}
    \nu \;\equiv\; -\,\frac{\varepsilon_{xx}}{\varepsilon_{zz}}
         \;=\;      -\,\frac{\varepsilon_{yy}}{\varepsilon_{zz}},
    \label{eq:lec02:poisson_def}
\end{align}
so that $\varepsilon_{xx} = \varepsilon_{yy} = -\nu\,\varepsilon_{zz}$.
The negative sign ensures $\nu > 0$ for any ordinary material (which
grows thinner when stretched).  As we shall show later from thermodynamic
stability arguments, $\nu$ is constrained to the interval
\[
    -1 < \nu \leq \tfrac{1}{2}.
\]
For virtually all engineering materials $\nu > 0$, i.e.\ the cross-section
shrinks under tension.

\noindent\textbf{Compression:} Reversing $F \to -F$ flips the sign of
all strains simultaneously; the formulas remain valid with the
understanding that $\sigma < 0$ denotes compression.

\subsubsection*{Typical values of elastic constants}

\begin{center}
\begin{tabular}{lccc}
\hline
\textbf{Material} & \textbf{$E$ (GPa)} & \textbf{$\nu$} & \textbf{Notes}\\
\hline
Steel (structural)  & 200     & 0.28 & Fe alloy, most common structural metal\\
Diamond             & 1\,000  & 0.20 & Hardest known material\\
Concrete            & 20--30  & 0.20 & Weak in tension; used in compression\\
Aluminium           & 70      & 0.33 & Light structural metal\\
Gold                & 78      & 0.44 & Dense, ductile\\
Bone (cortical)     & 15--25  & 0.30 & Varies along grain direction\\
Wood (along grain)  & 10--15  & 0.35 & Strongly anisotropic\\
Rubber              & $\sim 0.01$ & $\approx 0.50$ & Nearly incompressible\\
Cork                & $\sim 0.05$ & $\approx 0.00$ & Used as bottle stoppers\\
\hline
\end{tabular}
\end{center}

\textbf{Key observations:}
\begin{itemize}
    \item $E$ spans four orders of magnitude across common materials.
    \item $\nu$ is restricted to a much narrower range; most materials
          cluster between 0.25 and 0.45.
    \item Cork's near-zero $\nu$ is precisely what makes it an ideal
          bottle stopper: it does not expand laterally under axial
          compression, so it can be pushed into a bottle without cracking
          the glass.
    \item Rubber's $\nu \approx \frac{1}{2}$ means it is nearly
          \emph{incompressible}: volume is essentially conserved under
          deformation.
\end{itemize}

%%────────────────────────────────────────────────────────────────────
\subsection{Worked Example — Steel Column Under Compression}
\label{subsec:lec02:example}
%%────────────────────────────────────────────────────────────────────

\ex{}{\textbf{Problem.}
A structural steel column has length $L = 3\;\text{m}$, a square
cross-section of side $10\;\text{cm}$, and is subjected to a compressive
load $F = 200\;\text{kN}$.
\begin{itemize}
    \item For steel: $E = 200\;\text{GPa}$, $\nu = 0.28 \approx 0.3$.
    \item Find: the axial shortening $|\Delta L|$ and the lateral
          expansion $|\Delta x|$.
\end{itemize}

\medskip
\textbf{Step 1 — Convert units carefully.}
\[
    A = (0.10\;\text{m})^2 = 0.01\;\text{m}^2, \quad
    F = 200 \times 10^3\;\text{N}.
\]

\textbf{Step 2 — Compute the stress.}
\[
    \sigma = \frac{F}{A} = \frac{2 \times 10^5\;\text{N}}{10^{-2}\;\text{m}^2}
           = 2 \times 10^7\;\text{Pa} = 20\;\text{MPa}.
\]
(Note: the load is compressive, so strictly $\sigma = -20\;\text{MPa}$.)

\textbf{Step 3 — Compute the axial strain.}
\[
    \varepsilon_{zz} = \frac{\sigma}{E}
        = \frac{2 \times 10^7\;\text{Pa}}{2 \times 10^{11}\;\text{Pa}}
        = 10^{-4}.
\]

\textbf{Step 4 — Find $\Delta L$.}
\[
    |\Delta L| = |\varepsilon_{zz}|\,L = 10^{-4} \times 3\;\text{m}
              = 3 \times 10^{-4}\;\text{m}
              \;=\; \boxed{0.3\;\text{mm}}.
\]

\textbf{Step 5 — Find the lateral expansion.}
\[
    |\varepsilon_{xx}| = \nu\,|\varepsilon_{zz}|
        = 0.3 \times 10^{-4} = 3\times 10^{-5}.
\]
\[
    |\Delta x| = |\varepsilon_{xx}|\,x_0
        = 3 \times 10^{-5} \times 0.10\;\text{m}
        = 3 \times 10^{-6}\;\text{m}
        \;=\; \boxed{3\;\mu\text{m}}.
\]

\textbf{Sanity check.}  The column shortens by 0.3 mm (1 part in
$10^4$ of its length) and its cross-section widens by only 3 $\mu$m.
These tiny deformations justify the linear elastic (small-strain)
approximation: the geometry is essentially unchanged during loading.
The model breaks down when $|\varepsilon| \gtrsim 10^{-3}$, where
nonlinear or plastic effects become significant.

\textbf{Common student error:} Confusing $1\;\text{GPa} = 10^9\;\text{Pa}$
with $1\;\text{MPa} = 10^6\;\text{Pa}$ leads to an error of a factor of
$10^3$ in strain — just as was discussed in class.}

\begin{center}
\begin{tikzpicture}
  \def\L{5.0}
  \def\Rx{0.25}
  \def\Ry{0.5}
  \def\F{1.0}
  \def\x{0.14*\L}
  \draw[dashed]
    %(0,\Ry) arc(90:270:{\Rx} and {\Ry})
    %(0,0) ellipse ({\Rx} and {\Ry})
    (\x,\Ry) -- (0,\Ry) arc(90:270:{\Rx} and {\Ry})
    (0,-\Ry) -- (\x,-\Ry);
  \draw[<->] (0,-1.5*\Ry) --++ (\L,0) node[midway,above=-5,fill=white,inner sep=1] {$L$};
  \draw[width] (0, 1.3*\Ry) --++ (\x,0) node[midway,above] {$\Delta L$};
  \draw[force] (\L+\F,0) --++ (-\F,0) node[pos=0.4,above=-1] {$\vb{F}$};
  \draw[metal]
    (\x,-\Ry) |- (\L,\Ry) arc(90:-90:{\Rx} and {\Ry}) -- cycle;
  \draw[metal] (\x,0) ellipse ({\Rx} and {\Ry});
  \draw[force] (\x-\F,0) --++ (\F,0) node[pos=0,above=1,left=-2] {$\vb{F}$};
  \draw[dashed] (0,\Ry) arc(90:-90:{\Rx} and {\Ry});
\end{tikzpicture}
\end{center}

%%────────────────────────────────────────────────────────────────────
\subsection{The Stress Tensor and Its Symmetry}
\label{subsec:lec02:tensor}
%%────────────────────────────────────────────────────────────────────

\subsubsection*{From 1D to 3D: review of $\sigma_{ij}$}

The 1D example generalises: the force $d\mathbf{F}$ on any oriented
surface element $d\mathbf{S} = \hat{n}\,dA$ inside a solid is given by
the \textbf{stress tensor} $\bm{\sigma}$,
\begin{align}
    dF_i = \sigma_{ij}\,dS_j,
    \label{eq:lec02:cauchy_traction}
\end{align}
where summation over the repeated index $j$ is implied (Einstein
convention).  The tensor $\bm{\sigma}$ has $3\times 3 = 9$ components.
The diagonal components $\sigma_{11},\sigma_{22},\sigma_{33}$ are
\emph{normal stresses}; the off-diagonal components $\sigma_{ij}$ ($i\ne j$)
are \emph{shear stresses}.

\textbf{Stress vs.\ pressure:} Pressure in a fluid is isotropic —
$\sigma_{ij} = -P\delta_{ij}$ — and acts equally in all directions.
In a solid, the stress tensor need not be isotropic: different faces of
the same volume element may carry different forces.  Stress is the more
general concept.

\begin{center}
\tdplotsetmaincoords{70}{110}
\begin{tikzpicture}[tdplot_main_coords,line cap=round,line join=round,
    thick,>=Stealth,pics/dreibein/.style={code={%
      \foreach \XX [count=\YY from 0] in {#1}
      {
      \draw[->] (0,0,0) -- ($\XX1*(\dreibein\YY)$)
        node[above,dreibein/label]{$\sigma_{\pgfkeysvalueof{/tikz/dreibein/d}
        \the\numexpr\YY+1}$};}
    }},dreibein/.cd,label/.style={},d/.initial={}]
 \draw[->] (0,0,0) coordinate (O) -- (4,0,0) node[above]{$x_1$};
 \draw[->]  (O) -- (0,6,0) node[above]{$x_2$};
 \draw[->]  (O) -- (0,0,6) node[left]{$x_3$};
 \def\dreibein#1{\ifcase#1
 {1,0,0}\or{0,1,0}\or{0,0,1}\fi}
 \def\lstcols{"red","yellow","orange"}
 \begin{scope}[shift={(0,4,4)}]
  \foreach \X [evaluate=\X as \Y using {int(Mod(\X+1,3))},
    evaluate=\X as \Z using {int(Mod(\X+2,3))}]in {0,1,2}
  {\path ($-2*(\dreibein\X)$) pic[dreibein/label/.style={opacity=0}]{dreibein={-,-,-}};
   \pgfmathsetmacro{\mycolor}{{\lstcols}[\X]}
   \draw[fill=\mycolor,fill opacity=0.2] ($-2*(\dreibein\X)+2*(\dreibein\Y)+2*(\dreibein\Z)$)
    -- ($-2*(\dreibein\X)+2*(\dreibein\Y)-2*(\dreibein\Z)$)
    -- ($-2*(\dreibein\X)-2*(\dreibein\Y)-2*(\dreibein\Z)$)
    -- ($-2*(\dreibein\X)-2*(\dreibein\Y)+2*(\dreibein\Z)$) -- cycle;
   }
  \foreach \X [evaluate=\X as \Y using {int(Mod(\X+1,3))},
    evaluate=\X as \Z using {int(Mod(\X+2,3))}]in {0,1,2}
  {\pgfmathsetmacro{\mycolor}{{\lstcols}[\X]}
    \draw[fill=\mycolor,fill opacity=0.4] ($2*(\dreibein\X)+2*(\dreibein\Y)+2*(\dreibein\Z)$)
    -- ($2*(\dreibein\X)+2*(\dreibein\Y)-2*(\dreibein\Z)$)
    -- ($2*(\dreibein\X)-2*(\dreibein\Y)-2*(\dreibein\Z)$)
    -- ($2*(\dreibein\X)-2*(\dreibein\Y)+2*(\dreibein\Z)$) -- cycle;
   \path ($2*(\dreibein\X)$) pic[dreibein/d=\the\numexpr\X+1]{dreibein={+,+,+}};
   }
 \end{scope}
\end{tikzpicture}
\end{center}


\subsubsection*{Symmetry of $\sigma_{ij}$: torque argument}

We now show that $\sigma_{ij} = \sigma_{ji}$ must hold for any volume
element in static equilibrium.

Consider a small rectangular element with sides $\Delta x$, $\Delta y$,
$\Delta z$.  Focus on the torque about the $z$-axis.  The shear stress
$\sigma_{12}$ (the $y$-component of the force on the $+x$-face) exerts
a clockwise torque:
\[
    \tau_z^{(1)} = \sigma_{12}\,(\Delta y\,\Delta z)\cdot\frac{\Delta x}{2}
                  + \sigma_{12}\,(\Delta y\,\Delta z)\cdot\frac{\Delta x}{2}
                 = \sigma_{12}\,\Delta x\,\Delta y\,\Delta z.
\]
The corresponding shear $\sigma_{21}$ (the $x$-component of the force on
the $+y$-face) exerts a counter-clockwise torque of equal magnitude
$\sigma_{21}\,\Delta x\,\Delta y\,\Delta z$.  Rotational equilibrium
($\tau_z = 0$) requires:
\begin{align}
    \sigma_{12} = \sigma_{21}.
    \label{eq:lec02:sigma_sym_12}
\end{align}
Repeating for the other two planes gives $\sigma_{13}=\sigma_{31}$ and
$\sigma_{23}=\sigma_{32}$.  Hence:
\begin{align}
    \boxed{\sigma_{ij} = \sigma_{ji}},
    \label{eq:lec02:sigma_sym}
\end{align}
i.e.\ $\bm{\sigma}$ is a \textbf{symmetric} tensor with only
\textbf{6 independent components}.

\subsubsection*{Symmetry of $\sigma_{ij}$: Landau--Lifshitz argument}

An equivalent, more algebraic proof uses the total torque $\mathbf{N}$
on a volume element:
\[
    N_i \;=\; \oint_{\partial V} \varepsilon_{ijk}\,x_j\,\sigma_{kl}\,dS_l
         \;=\; \int_V \varepsilon_{ijk}\,\partial_l(x_j\,\sigma_{kl})\,dV.
\]
Applying the product rule, $\partial_l(x_j\sigma_{kl})
= \delta_{jl}\,\sigma_{kl} + x_j\,\partial_l\sigma_{kl}$.
The second term is the moment of the body-force density $f_k = \partial_l\sigma_{kl}$,
which represents the torque due to net forces on the element and is
already accounted for in the force-balance equation.  The first term
must vanish independently:
\[
    \varepsilon_{ijk}\,\delta_{jl}\,\sigma_{kl}
    = \varepsilon_{ijk}\,\sigma_{kj} = 0.
\]
For $i=3$: $\varepsilon_{312}\sigma_{21} + \varepsilon_{321}\sigma_{12}
    = \sigma_{21} - \sigma_{12} = 0$,
i.e.\ $\sigma_{12}=\sigma_{21}$.  The same follows for all pairs,
confirming~\eqref{eq:lec02:sigma_sym}.

%%────────────────────────────────────────────────────────────────────
\subsection{Static Equilibrium Equations}
\label{subsec:lec02:equil}
%%────────────────────────────────────────────────────────────────────

The net force on a volume element $dV$ comes from two sources:
\begin{enumerate}
    \item \textbf{Contact (internal) forces} from surrounding material,
          giving a force per unit volume $\partial_j\sigma_{ij}$
          (cf.\ the divergence theorem applied to
          Eq.~\eqref{eq:lec02:cauchy_traction}).
    \item \textbf{Body forces} — for gravity this is
          $\rho\,g_i$ per unit volume, where $\rho$ is mass density.
\end{enumerate}

Newton's second law for static equilibrium ($\mathbf{a}=0$) yields:
\begin{align}
    \boxed{\partial_j\,\sigma_{ij} + \rho\,g_i \;=\; 0}.
    \label{eq:lec02:equilibrium}
\end{align}

This is the Cauchy equilibrium equation.  It is the continuum analogue of
$\sum F = 0$ for a particle.  The gradient of stress (divergence of
$\bm{\sigma}$) must balance the body-force density.

\textbf{Boundary conditions:}  At a free surface with outward normal
$\hat{n}$, the traction must equal the applied surface load $t_i$:
\[
    \sigma_{ij}\,n_j \Big|_{\text{surface}} = t_i.
\]
For a truly free surface (no applied load), $t_i = 0$.

\textbf{Remark:}  The precise sign convention for Eq.~\eqref{eq:lec02:equilibrium}
and the boundary conditions will be refined at the start of Lecture~03.

%%────────────────────────────────────────────────────────────────────
\subsection{Brief Note on Tensors}
\label{subsec:lec02:tensors}
%%────────────────────────────────────────────────────────────────────

A \textbf{tensor of rank $n$} is a multilinear map that takes $n$ vectors
and returns a scalar, represented in a given basis by an array of
$3^n$ components.  Familiar examples:
\begin{itemize}
    \item Rank-0 (scalar): temperature $T$.
    \item Rank-1 (vector): displacement $u_i$, body force $g_i$.
    \item Rank-2: stress $\sigma_{ij}$, strain $\varepsilon_{ij}$.
    \item Rank-4: the elastic stiffness tensor $C_{ijkl}$  
          (to be introduced shortly) relates stress and strain in 3D.
\end{itemize}

%%────────────────────────────────────────────────────────────────────
\subsection*{Key Takeaways}
\label{subsec:lec02:takeaways}
%%────────────────────────────────────────────────────────────────────

\begin{itemize}
    \item \textbf{Stress} $\sigma = F/A$ is force per unit area
          (units: Pa).  It is the quantity that drives deformation, not
          the bare force.
    \item \textbf{Longitudinal strain} $\varepsilon_{zz} = \Delta L / L$
          is the dimensionless relative elongation.
    \item \textbf{Hooke's law} in 1D: $\sigma = E\,\varepsilon_{zz}$,
          giving $\Delta L = FL / (EA)$.  Linear elastic response holds
          for $|\varepsilon| \lesssim 10^{-3}$.
    \item \textbf{Young's modulus} $E$ (units: Pa) characterises the
          stiffness of a material.  It spans from $\sim 0.01\;\text{GPa}$
          (rubber) to $\sim 1000\;\text{GPa}$ (diamond).
    \item \textbf{Poisson's ratio} $\nu$ (dimensionless) relates lateral
          contraction to longitudinal extension:
          $\varepsilon_{xx} = -\nu\,\varepsilon_{zz}$.
          Physically $0 < \nu \leq \tfrac{1}{2}$ for most materials;
          thermodynamics requires $-1<\nu\leq\tfrac{1}{2}$.
    \item The \textbf{stress tensor} $\sigma_{ij}$ generalises 1D stress
          to 3D.  It is \textbf{symmetric} ($\sigma_{ij}=\sigma_{ji}$)
          because rotational equilibrium must hold for every volume element,
          leaving only \textbf{6 independent components}.
    \item \textbf{Equilibrium equation}:
          $\partial_j\sigma_{ij} + \rho g_i = 0$, the continuum version
          of $\sum F_i = 0$.
\end{itemize}

